\documentclass[12pt]{article}
\usepackage[utf8]{inputenc}
\usepackage[T1]{fontenc}
\usepackage{hyperref}
\title{Tauya Language Resources for K-12 Curriculum Development}
\author{Compiled by Claude Sonnet 5 (claude-sonnet-5) for Andr\'as Kornai}
\date{}
\begin{document}
\maketitle

\section{Introduction}
Tauya (ISO 639-3 tya) is a Rai Coast language of Madang Province, Papua New Guinea, with about 350 speakers. Most of what exists is linguistic description by Lorna MacDonald.

\section{Primary Educational Resources}
No primers, textbooks or school materials were found.

\section{Linguistic and Reference Materials}
\begin{itemize}
\item MacDonald, \emph{A Grammar of Tauya}, thesis (1985), open access, University of Manitoba: \url{https://mspace.lib.umanitoba.ca/items/8283d204-6b0b-4cb6-9339-153296e37d66/full}
\item MacDonald, \emph{A Grammar of Tauya}, Mouton Grammar Library 6 (1990), catalogue record: \url{https://www.oeaw.ac.at/resources/Record/9783110846027?lng=de}
\item Review of the grammar by Carl R. Whitehead (1993): \url{https://glottolog.org/resource/reference/id/52304}
\item MacDonald, \emph{A Dictionary of Tauya} (Pacific Linguistics 638, 2013, over 2,500 entries); this was reported in a search summary and no direct page was seen.
\item Wikipedia article on Tauya: \url{https://Www.Wikipedia.org/wiki/Tauya}
\item WALS language page: \url{https://wals.info/languoid/lect/wals_code_tau}
\end{itemize}

\section{Religious and Community Resources}
None found.

\section{Digital and Audio Resources}
None found.

\section{Curriculum Development Context}
The grammar and dictionary are a basis for writing an orthography and teaching materials, but no existing school materials were found.
\end{document}